%%==========================================================
%% CAPÍTULO 3: ANÁLISIS DEL PROBLEMA Y MARCO DE REFERENCIA DE LA SOLUCIÓN
%%==========================================================
\label{cap:fundamentacion}

La primera fase de la metodología del proyecto (Figura~\ref{fig:metodologia})
persiguió dos propósitos, el primero de los cuales fue demostrar con evidencia
que el problema existe y establecer en qué consiste exactamente, y el segundo
definir cómo debía estructurarse la solución antes de escribir la primera
línea de código. Al cerrarla quedó establecido qué debe hacer la herramienta y
qué forma deben tener los proyectos que produzca.

El trabajo se organizó en dos bloques, y el primero reunió la
evidencia que sustenta el proyecto, levantando el contexto donde la
herramienta iba a operar, el testimonio de quienes trabajan en él y lo que
ofrecen las soluciones existentes, hasta especificar formalmente qué debía
construirse. El segundo definió el encuadre de la solución y los fundamentos
sobre los que se apoya, que es diseño de alto nivel y llega hasta ahí, porque
las reglas concretas que gobiernan la forma de lo que la línea produce se
definen en la fase que las construye y las verifica.

Lo que este capítulo presenta es una síntesis del proceso y del argumento que
de él resultó, mientras que el detalle completo de cada instrumento, las
transcripciones de las entrevistas, la evaluación de cada solución frente a
cada eje y la ficha de cada requerimiento se conserva en los documentos de la
Fase 1, disponibles en la carpeta \texttt{Implementación/} del repositorio del
trabajo de
grado\footnote{\url{https://github.com/facundo1220/BackendServiceLine}}.

\section{Comprensión del problema y especificación}
\label{sec:comprension-problema}

La dimensión del problema de la fase consistió en aplicar un proceso de
obtención de requisitos sobre un caso de estudio, para comprender cómo se
inicia un proyecto backend y cómo se integran funcionalidades reutilizables en
la práctica real de un equipo de desarrollo.
El caso de estudio fue el Laboratorio Digital del Banco W
\citep{UniversidadIcesi2023}, cuyas prácticas, roles y restricciones
alimentaron la especificación de requerimientos de la herramienta.

\subsection{Caracterización del contexto de aplicación}
\label{sec:caracterizacion}

El Laboratorio Digital opera bajo marcos de trabajo ágiles, con Scrum y una
práctica de Dual Track que alinea el discovery con el delivery en sprints de
dos semanas \citep{MartyCagan2012}, y su composición de roles corresponde a la
de un equipo ágil, con desarrolladores, testers, un líder técnico y líderes de
producto. Su equipo técnico es reducido y estable, con cuatro desarrolladores
que atienden varios productos a la vez.

Como todo equipo que construye software dentro de una organización regulada,
trabaja en dos planos a la vez, porque conserva autonomía sobre la
construcción interna de sus soluciones, que es donde decide cómo se estructura
y se escribe el código de cada producto, y al mismo tiempo se subordina a las
definiciones de las funciones institucionales de tecnología, que establecen
los estándares a los que todo desarrollo debe ajustarse y las condiciones bajo
las cuales un servicio se integra con los sistemas centrales y llega a operar.
Esa frontera delimita sobre qué puede actuar legítimamente una herramienta de
generación, porque lo que produce pertenece por entero al primer plano y nada
de lo que genera puede alterar lo que la organización impone.

\subsection{Identificación y justificación de stakeholders}
\label{sec:stakeholders}

Definido el contexto, el paso siguiente fue establecer a quiénes consultar, y
la selección de roles condiciona lo que se va a encontrar, porque un
arquitecto, un desarrollador y un líder técnico ven cosas distintas del mismo
problema, y por eso el criterio se fijó antes de entrevistar. Se consideró
stakeholder a quien usa la herramienta, a quien fija reglas que esta debe
respetar, a quien puede habilitar o bloquear su adopción, y a quien ve cambiar
su trabajo por ella aunque no la use. Para ordenarlos se adoptó el modelo de
capas (\emph{onion model}) de Alexander \citep{Alexander2005}, que los agrupa
según su cercanía al artefacto, desde quienes lo operan hasta quienes solo
imponen condiciones sobre él.

La Tabla~\ref{tab:roles} recoge los roles consultados, la capa a la que
pertenecen y lo que se esperaba de la perspectiva de cada uno.

\begin{table}[H]
\centering
{\footnotesize
\begin{tabular}{p{3.2cm}p{2.6cm}p{9.4cm}}
\toprule
\textbf{Rol} & \textbf{Capa} & \textbf{Aporte esperado} \\
\midrule
Desarrolladores backend & Núcleo & Desagregación del proceso de arranque, puntos de dolor concretos y criterio sobre qué estandarizar \\
Testers & Núcleo & Requerimientos de consistencia externa de APIs e interfaces, y evidencia de heterogeneidad entre productos \\
Líder técnico & Sistema contenedor & Visión consolidada de la divergencia entre proyectos y consideraciones de gobernanza del catálogo \\
Líder de producto & Sistema contenedor & Impacto del tiempo de arranque e integración sobre la planeación de productos \\
Arquitectura & Entorno & Condiciones de diseño de un módulo de integración, organización de dominios y criterios de evolución \\
Desarrolladores externos & Entorno & Contraste sobre qué problemas son generales de la industria y qué mecanismos de reutilización aplican otras organizaciones \\
\bottomrule
\end{tabular}
}
\caption{Roles consultados, por capa del modelo de Alexander}
\label{tab:roles}
\end{table}

Los cuatro primeros roles pertenecen al equipo estudiado y los dos últimos no.
La función de arquitectura representa las definiciones a las que cualquier
desarrollo se subordina en una organización regulada, y fue la única fuente
que aportó condiciones de diseño en lugar de descripción de la práctica. Los
desarrolladores externos permiten separar lo propio del contexto estudiado de
lo que ocurre también en otras organizaciones. Las funciones consultadas son
arquetipos de cualquier equipo de producto con estas características, y es eso
lo que permite trasladar el análisis a otro contexto del mismo tipo.

\subsection{Protocolo de recolección de información}
\label{sec:protocolo}

La información se recolectó mediante entrevistas semiestructuradas de diseño
mixto, con nueve entrevistas repartidas en seis guiones, uno por cada rol
consultado. Una primera parte de preguntas abiertas recogió el proceso actual,
los puntos de dolor y el criterio de cada entrevistado, y es la que sostiene el
análisis cualitativo del que se derivan los temas, mientras que una segunda
parte más breve, de preguntas cerradas y respondida en los minutos finales,
produjo datos agregables que complementan las respuestas abiertas.

Esas preguntas cerradas se aplicaron antes de que existiera la primera línea de
código de la herramienta, de modo que son la única fuente de autoreporte
anterior al artefacto, y sobre ellas se constituyó la línea base con la que
compara la evaluación del Capítulo~\ref{cap:evaluacion}.

\subsection{Análisis cualitativo y temas derivados}
\label{sec:analisis-cualitativo}

Al ser las entrevistas de preguntas abiertas, cada entrevistado describió su
propia experiencia con sus propias palabras, de modo que el material se
transcribió y se analizó buscando qué se repetía de un testimonio a otro,
con el fin de encontrar los patrones de la problemática y no opiniones
aisladas. De ese análisis salieron quince temas, que recogen lo que los
entrevistados describieron de su práctica sin redactarlo como requerimiento ni
proponer una solución, y que se agruparon después en los cinco conjuntos de la
Tabla~\ref{tab:grupos}, porque es la agrupación la que permite analizarlos y la
que da origen a los requerimientos.

\begin{table}[H]
\centering
{\footnotesize
\begin{tabular}{p{2.9cm}p{1.7cm}p{6.9cm}p{3.1cm}}
\toprule
\textbf{Grupo} & \textbf{Temas} & \textbf{Lo que el testimonio muestra} & \textbf{Roles que lo evidencian} \\
\midrule
El costo del arranque & T-01, T-02, T-03 & Arrancar un proyecto es un proceso manual, concentrado en una persona, que impone un costo cognitivo inicial, se paga en alcance de negocio y resulta invisible para el negocio & Desarrolladores, líder técnico, líder de producto, externos \\
La reutilización y su propagación & T-04, T-05, T-09 & Existe un núcleo reconocido de funcionalidades que se repite entre productos, pero sin un mecanismo formal se recurre a copiarlo, y el contrato del servicio es el punto más frágil al propagar un cambio & Desarrolladores, arquitectura, líder de producto, externos \\
Causas estructurales de la divergencia & T-06, T-07, T-08 & Los proyectos divergen porque nada propaga las mejoras ni refuerza el estándar, porque se diseña para el alcance inmediato y porque el conocimiento transversal llega tarde & Líder de producto, líder técnico, arquitectura, externos \\
Condiciones de diseño & T-10, T-11, T-12, T-13 & La solución debe estandarizar con rigidez justificada, separar lo fijo de lo variable, evolucionar por versiones y poder verificarse & Todos los roles consultados \\
Condiciones de adopción y permanencia & T-14, T-15 & Adoptar requiere agilidad real y uso sencillo; permanecer requiere capacidad de evolución y un responsable definido & Desarrolladores, líder técnico, arquitectura, externos \\
\bottomrule
\end{tabular}
}
\caption{Los quince temas del análisis cualitativo, por grupo}
\label{tab:grupos}
\end{table}

El respaldo de los temas no es uniforme, porque catorce de los quince los
evidencia más de un rol, y tres de ellos aparecen en cuatro o cinco roles
distintos, lo que los sitúa como hallazgos compartidos entre perspectivas que
miran el mismo problema desde lugares diferentes. El único que sostiene un
solo rol es el que afirma que lo fijado en el código impide reutilizarlo
después, evidenciado por el líder de producto, y se conservó porque introduce
una causa estructural que ningún otro nombró, con el respaldo menor que eso
supone.

Los quince temas, agrupados así, son la unidad de trabajo del resto del
proyecto. Todo requerimiento posterior cita el tema o los temas que lo
originan, y el cierre del proyecto, en el Capítulo~\ref{cap:evaluacion}, se
formula contra los requerimientos y las brechas.

\subsection{Estudio comparativo de soluciones existentes}
\label{sec:estudio-comparativo}

Establecido qué muestra cada grupo temático, el paso siguiente fue medir qué
parte de eso resuelven ya las soluciones disponibles y qué parte no resuelve
ninguna, que es lo que determina si el problema planteado en la
sección~\ref{sec:planteamiento} queda cubierto por una herramienta existente.
La comparación se hizo sobre las ocho soluciones descritas en el estado del
arte (sección~\ref{sec:estado-del-arte}), que el Capítulo~\ref{cap:marco}
describe una por una y que esta sección evalúa frente a los pilares que los
temas establecieron, sin volver a describir ninguna.

Cada eje del marco de evaluación responde a un tema documentado con evidencia
de entrevistas, en lugar de a un inventario genérico de funcionalidades, de
modo que la comparación mide la distancia entre cada solución y el problema
real del contexto. Se excluyó el tema T-03, por tratarse de un efecto
organizacional y no de una capacidad evaluable en una herramienta, y los
catorce restantes quedan cubiertos por al menos un eje de la
Tabla~\ref{tab:ejes}, sin que ningún eje quede sin tema de origen.

\begin{table}[H]
\centering
{\footnotesize
\begin{tabular}{p{3.2cm}p{5.3cm}p{1.2cm}p{4.6cm}}
\toprule
\textbf{Enfoque} & \textbf{Eje} & \textbf{Temas} & \textbf{Qué significa cumplirlo} \\
\midrule
Generación de proyectos & Generación desde plantilla & T-01, T-02 & Producir un proyecto ejecutable a partir de una plantilla parametrizable, sin intervención manual \\
 & Contenido resuelto de origen & T-07, T-08 & Que el proyecto traiga ya resueltas las capacidades transversales y las restricciones que suelen descubrirse tarde \\
 & Externalización de la configuración & T-11 & Que lo susceptible de cambiar quede fuera del código y lo estable no \\
\midrule
Ciclo de vida & Composición de módulos en proyectos existentes & T-04, T-09 & Incorporar una capacidad transversal a un proyecto ya en curso, de forma desacoplada \\
 & Propagación de cambios de una capacidad compartida & T-05 & Que un cambio en una capacidad alcance a los proyectos que la usan sin intervenirlos uno por uno \\
 & Consistencia entre proyectos en el tiempo & T-06 & Detectar y corregir la desviación entre proyectos generados con versiones distintas \\
 & Versionado y distribución & T-12 & Que las capacidades y la herramienta evolucionen por versiones con adopción gradual \\
\midrule
Adopción y gobernanza & Adaptabilidad de lo generado & T-10 & Que el resultado se pueda modificar sin romper la relación con la herramienta \\
 & Verificabilidad & T-13 & Que lo generado y la herramienta puedan someterse a verificación automática \\
 & Costo de uso & T-14 & Que generar un proyecto cueste un comando o pocos comandos, sin configuración previa extensa \\
 & Gobernanza del catálogo & T-15 & Que exista un catálogo con propiedad, control de acceso y reglas de incorporación \\
\bottomrule
\end{tabular}
}
\caption{Los once ejes del marco de evaluación y el tema que cada uno atiende}
\label{tab:ejes}
\end{table}

La evaluación de las ocho soluciones frente a los once ejes, celda por celda,
se conserva en el documento \emph{Estudio comparativo de soluciones
existentes} de la Fase 1, disponible en el repositorio del trabajo de grado.
Su lectura agregada es que ninguna cubre los tres enfoques a la vez. Las
herramientas de plantilla resuelven la generación y el costo de uso, pero no
la composición sobre un proyecto en curso ni la gobernanza; las que sostienen
el ciclo de vida toman la plantilla completa como unidad de cambio; y las
plataformas de desarrollo interno resuelven el catálogo y la propiedad al
precio de desplegar y operar una aplicación propia.

\paragraph*{Brechas identificadas}
\label{sec:brechas}

Del cruce entre las ocho soluciones y los once ejes surgen tres capacidades
que el contexto necesita y que ninguna herramienta del mercado resuelve.

\textbf{B-01. Incorporar un módulo a un proyecto en curso.} Cuando meses
después un proyecto necesita autenticación, observabilidad o el consumo de un
servicio compartido, no hay forma de incorporarlo salvo a mano. La solución
más cercana modifica código existente en lugar de limitarse a crear archivos
nuevos, pero exige que todos los proyectos convivan en un mismo repositorio,
y el contexto estudiado trabaja con repositorios independientes.

\textbf{B-02. Propagar el cambio de un módulo.} Cuando una capacidad
compartida se copió a varios proyectos y después cambia, el cambio debe
llegar a todos sin intervenir cada uno a mano. Las soluciones más cercanas
registran la versión de plantilla que usó el proyecto y aplican los cambios
preservando las modificaciones locales, pero su unidad es la plantilla
completa. Lo que cambia en el caso documentado es una capacidad puntual y no
la estructura, de modo que actualizar la plantilla entera no propaga el
cambio de un módulo.

\textbf{B-03. Catálogo con reglas y bajo costo de uso.} Sostener un catálogo
exige propiedad definida, control de acceso y reglas de incorporación, sin
que eso encarezca el uso diario. Las soluciones más cercanas lo ofrecen, pero
exigen desplegar, configurar y operar una aplicación propia antes de generar
el primer proyecto, un costo desproporcionado para un equipo de cuatro
desarrolladores. En el extremo opuesto, las herramientas de un comando no
tienen catálogo alguno.

Las tres brechas cumplen una doble función, porque justifican construir algo
nuevo en lugar de adoptar una herramienta existente, y funcionan como criterio
de prioridad, porque un requerimiento que cierra una brecha es imprescindible
mientras que otro que solo mejora la experiencia de uso puede diferirse.

\subsection{Especificación de requerimientos}
\label{sec:especificacion-requerimientos}

De los quince temas, que expresan qué espera resolver el equipo, y de las tres
brechas, que delimitan lo que ninguna solución existente cubre, salieron
veintidós requerimientos, dieciséis funcionales y seis no funcionales. Con
ellos cierra el bloque de evidencia, porque son la forma verificable de todo
lo que las subsecciones anteriores establecieron.

A diferencia del tema, que describe lo que ocurre en la práctica, el
requerimiento fija un alcance concreto y establece qué debe hacer la
herramienta, sobre qué actúa y bajo qué condición se da por satisfecho. Por
eso cada uno se acompaña de su criterio de aceptación, que es el que vuelve
verificable aquello que el tema solo enunciaba, y la ficha de los veintidós,
con su enunciado, su prioridad, su origen y su criterio, se conserva en el
documento \emph{Especificación de requerimientos} de la Fase 1, disponible en
el repositorio del trabajo de grado.

\subsubsection*{Prioridad y cobertura}

Los requerimientos se priorizaron con la escala MoSCoW, cuyos niveles se
asignan según el origen de cada uno y no según la preferencia del
implementador. Es imprescindible el que responde a una de las tres brechas o
sin el cual la solución no resuelve el problema levantado; importante, el que
responde a un tema documentado pero cuya ausencia no invalida la
solución en esta fase; y deseable, el que aporta valor comprobado y puede
quedar fuera si el alcance lo exige.

Los dieciséis requerimientos funcionales se enuncian sobre las dos capacidades
centrales de la herramienta, generar proyectos e incorporar módulos, y sobre
las dos condiciones que las sostienen en el tiempo, la actualización y el
catálogo. La Tabla~\ref{tab:rf} recoge qué exige cada una y de qué temas y
brechas proviene.

\begin{table}[H]
\centering
{\footnotesize
\begin{tabular}{p{3.1cm}p{7.6cm}p{1.6cm}p{2.3cm}}
\toprule
\textbf{Capacidad} & \textbf{Qué exige} & \textbf{Req.} & \textbf{Origen} \\
\midrule
Generación de proyectos & Producir un proyecto ejecutable y parametrizable, con la estructura y el entorno de trabajo del marco ya resueltos, su documentación de API automática y un resultado determinista ante la misma entrada & RF-01 a RF-05 & T-01, T-02, T-06, T-07, T-08, T-11 \\
Composición de módulos & Incorporar un módulo a un proyecto existente sin que nadie edite archivos a mano, modificando lo que deba referenciarlo sin sobrescribir código propio, con consulta del estado y al menos un módulo real en el catálogo & RF-06 a RF-09 & T-04, T-09, T-15, B-01 \\
Actualización y propagación & Registrar en el proyecto qué módulos incorporó y en qué versión, actualizarlos sin intervención manual, informar de los desactualizados y preservar el código del equipo al hacerlo & RF-10 a RF-13 & T-05, T-06, T-10, B-02 \\
Catálogo y distribución & Declarar el catálogo de forma declarativa y versionada sin componentes de servidor, distribuir por el registro de paquetes de la organización y registrar la propiedad de cada activo & RF-14 a RF-16 & T-12, T-15, B-03 \\
\bottomrule
\end{tabular}
}
\caption{Los dieciséis requerimientos funcionales, por capacidad de la herramienta}
\label{tab:rf}
\end{table}

Los seis requerimientos no funcionales de la Tabla~\ref{tab:rnf} expresan los
atributos de calidad que la herramienta y lo que produce deben satisfacer,
con independencia de las operaciones concretas que ofrezca.

\begin{table}[H]
\centering
{\footnotesize
\begin{tabular}{p{1.4cm}p{2.9cm}p{8.9cm}p{1.3cm}}
\toprule
\textbf{ID} & \textbf{Atributo} & \textbf{Requerimiento} & \textbf{Origen} \\
\midrule
RNF-01 & Usabilidad & Las operaciones de la herramienta deben ejecutarse con un comando cada una & T-14 \\
RNF-02 & Adaptabilidad & El código generado debe ser modificable sin que ello impida seguir usando la herramienta sobre el proyecto & T-10 \\
RNF-03 & Verificabilidad & Lo generado debe poder someterse a análisis estático de código, incluidos aspectos de seguridad & T-13 \\
RNF-04 & Verificabilidad & La herramienta debe poder someterse a auditoría de dependencias y vulnerabilidades, con resultado consultable & T-13 \\
RNF-05 & Configurabilidad & El proyecto generado no debe contener valores fijos susceptibles de cambio, en particular direcciones de servicios y parámetros de negocio & T-11 \\
RNF-06 & Evolución & La herramienta y sus módulos deben evolucionar por versiones, admitiendo la convivencia de versiones distintas entre proyectos & T-12 \\
\bottomrule
\end{tabular}
}
\caption{Requerimientos no funcionales}
\label{tab:rnf}
\end{table}

Con la especificación queda cerrado el bloque de evidencia, y estos veintidós
requerimientos son el insumo del bloque siguiente, donde el marco de
referencia tendrá que habilitar cada uno o declarar a qué fase de
construcción se difiere. La comprobación de que cada uno quedó satisfecho se
realiza al cierre del proyecto, cuando existe el artefacto contra el cual
verificarla.

\section{Definición del marco de referencia}
\label{sec:marco-tecnico-solucion}

Establecida la evidencia del problema, el marco de referencia fija bajo qué
enfoque de ingeniería se construye la solución y sobre qué fundamentos se
apoya, de modo que las fases siguientes construyan dentro de límites ya
decididos en lugar de fijarlos durante la implementación. Lo que el marco
define no es la herramienta, sino la familia de proyectos que la herramienta
produce, porque es esa familia la que tiene que compartir una estructura para
que un mismo activo pueda servir a varios de sus productos.

Se desarrolla en dos partes, de las cuales la primera encuadra la
solución como una línea de producto de software, que es el enfoque que ordena
la construcción de los activos reutilizables y la de los productos que se
derivan de ellos, y declara con ese vocabulario qué productos forman la
familia y cómo se organiza su construcción. La segunda adopta los fundamentos
que dan forma al interior de cada producto y fija hasta dónde se toma cada
uno. El marco termina ahí, porque las reglas concretas que gobiernan la forma
de lo que la línea produce se definen y se verifican en la fase que las
construye, y se conservan en los documentos de esa fase.

\subsection{La solución como línea de producto de software}
\label{sec:caracterizacion-lps}

Una línea de producto de software es una familia de sistemas que comparten un
conjunto común y gestionado de características y que se construyen a partir de
activos comunes siguiendo un procedimiento establecido \citep{Clements2001}.
Lo que la separa de una colección de proyectos parecidos es que esas
similitudes están declaradas y gestionadas, de manera que un activo se
construye una vez y sirve a varios productos, en lugar de rehacerse en cada
uno. Encuadrar la solución bajo ese enfoque es la decisión de fondo del
trabajo, porque desplaza el problema desde elegir una herramienta hacia
definir las condiciones que gobiernan su uso, que es el plano en el que las
tres brechas pueden cerrarse para una familia entera de productos y no para un
proyecto a la vez.

\subsubsection*{Elementos constitutivos y pertinencia del enfoque}

El enfoque exige cuatro elementos constitutivos, y la ausencia de cualquiera
de ellos impide hablar de una línea de producto \citep{Clements2001}. Lo que
vuelve pertinente el encuadre no es que los cuatro puedan construirse, sino
que los cuatro ya estaban presentes en el contexto estudiado de forma
informal, cada uno con su evidencia en el análisis del problema. La
Tabla~\ref{tab:constitutivos} enfrenta lo que el enfoque exige con lo que el
contexto tenía y con lo que la línea construyó en su lugar.

\begin{table}[H]
\centering
{\footnotesize
\begin{tabular}{p{2.5cm}p{3.9cm}p{4.3cm}p{3.8cm}}
\toprule
\textbf{Elemento} & \textbf{Qué exige el enfoque} & \textbf{Cómo estaba en el contexto} & \textbf{Cómo se materializa en la línea} \\
\midrule
Alcance declarado & Declaración explícita de qué productos pertenecen a la familia & Una familia sin declarar, cuyos productos comparten propósito y estructura sin que nadie lo haya escrito & El alcance de la familia, con su condición mínima de pertenencia y sus exclusiones \\
Arquitectura común & Estructura compartida que hace posible que un mismo activo sirva a varios productos & Un patrón que el desarrollador debe dominar antes de ser productivo, transmitido imitando otros proyectos & Domain-Driven Design y arquitectura hexagonal, con una regla de dependencia verificable \\
Core assets & Activos reutilizables que incluyen arquitectura, contratos y modelos, y no solo código & Un núcleo reconocido de funcionalidades transversales que se repite entre productos y se copia & Las plantillas, los módulos, el \emph{contract package} y el catálogo \\
Mecanismo de derivación & Procedimiento que produce un producto a partir de los activos y de una selección de características & Copiar un proyecto y ajustarlo, que deriva sin trazabilidad ni vínculo con su origen & La herramienta de línea de comandos, construida en las fases posteriores \\
\bottomrule
\end{tabular}
}
\caption{Los cuatro elementos constitutivos de una línea de producto, frente a lo que el contexto ya tenía y a lo que la línea construyó}
\label{tab:constitutivos}
\end{table}

La tercera columna es la que decide, porque los cuatro elementos existían,
pero ninguno declarado ni gestionado, y esa es la diferencia que el enfoque
convierte en ventaja de ingeniería, porque lo que no está declarado no puede
compartirse de forma deliberada. La cuarta fila lo muestra en su forma más
clara, porque copiar un proyecto y ajustarlo es lo que la literatura denomina
\emph{clone-and-own}, una forma legítima de reutilización por su bajo costo
inicial, cuyo problema característico es la divergencia progresiva, porque
nada devuelve a las demás copias lo que se corrige en una y el costo crece con
el número de copias \citep{Dubinsky2013}. Es exactamente lo que el análisis
del problema encontró en campo, y lo que la primera brecha enuncia como
capacidad ausente.

De ahí se sigue el papel que le toca al artefacto, porque lo que el encuadre
aporta no es una herramienta más, sino las condiciones bajo las cuales se
emplean las que ya existen, y la que las fases siguientes construyen es el
mecanismo de derivación que esas condiciones exigen. Es además lo que vuelve
gestionados a los otros tres elementos, porque un alcance declarado y unos
activos comunes sin un procedimiento que los aplique vuelven a quedar en manos
de quien copia.

\subsubsection*{Organización general de la línea}

Los cuatro elementos describen qué debe existir, pero no cómo se articulan, y
para articularlos se adopta el modelo de referencia que cruza el espacio del
problema con el de la solución, y la ingeniería de dominio con la de
aplicación \citep{Apel2013}. Del cruce resultan cuatro cuadrantes, donde el
análisis del dominio declara el alcance de la familia y la variabilidad; la
implementación del dominio construye los activos comunes y la arquitectura de
referencia; el análisis de requerimientos de la aplicación es el momento en
que un desarrollador selecciona características para un producto concreto; y
la derivación del producto las ensambla. Los dos primeros corresponden a las
fases 1, 2 y 3 de este trabajo; los dos últimos ocurren cada vez que alguien
emplea la herramienta.

De los cinco flujos que conectan los cuadrantes, el que más condiciona el
diseño es el retorno de nuevos requerimientos hacia el dominio, porque es el
que permite que lo aprendido construyendo corrija el encuadre. El modelo
supone además un modelo de características formal sobre el cual se ejecuta la
selección, y este marco no lo define, de modo que la variabilidad de la línea
queda asignada a la Fase 2, que deberá formalizarla dentro de los límites que
el alcance de la familia y el catálogo de activos imponen.

\subsubsection*{Alcance de la familia de productos}

Pertenecen a la familia los servicios backend sobre Node.js y TypeScript que
responden por una sola capacidad de negocio, exponen al menos una API HTTP
documentada bajo el estándar OpenAPI y consumen servicios corporativos. Quedan
fuera las aplicaciones frontend, las que no exponen una API HTTP, los scripts
de propósito único y los servicios construidos sobre otros lenguajes o
ecosistemas.

Esa exposición es la condición mínima de pertenencia, y de ella depende que
todos los productos compartan la misma superficie de entrada, que es lo que
hace viable un activo común y lo que da sentido a la documentación OpenAPI
como capacidad transversal de la línea. Delimitar el alcance con esta
precisión lo exige el propio enfoque, porque uno demasiado amplio impide que
un activo sirva a todos los productos y uno demasiado estrecho no justifica el
esfuerzo de construir la línea frente al de construir un solo producto
\citep{Clements2002}.

La condición no fija un modelo de ejecución, porque un producto sin servidor
propio la cumple igual, y cuántas plantillas tiene la línea lo decide el
mecanismo de derivación y no el alcance de la familia. Que el marco absorbiera
ese hallazgo de la construcción sin alterar su condición de pertenencia es la
primera evidencia del flujo de retorno que el modelo de referencia declara.

\subsubsection*{Desviaciones declaradas respecto del modelo clásico}

El modelo adoptado opera sobre dos asunciones que el contexto no satisface, y
se declaran como desviaciones explícitas en lugar de forzar el caso dentro de
un esquema que no lo previó.

La primera es la de productos autónomos, porque el modelo no contempla retorno
desde el producto derivado, que una vez producido queda fuera del modelo,
mientras que aquí el proyecto se entrega a un equipo que conserva autonomía
para modificarlo. La línea no gobierna el producto después de derivado, sino
que mantiene con él un vínculo que permite propagar cambios sin imponer su
aceptación.

La segunda es la de derivación por etapas, porque en el modelo la selección de
características ocurre una sola vez y precede a la derivación, mientras que
aquí puede repetirse sobre un producto ya derivado, porque una capacidad
compartida debe poder incorporarse a un proyecto en curso.

Ambas se originan en la misma condición, que es que los productos pertenecen
a equipos autónomos con repositorios independientes. No son concesiones, sino
la consecuencia de aplicar el enfoque a un contexto que el modelo clásico no
previó.

\subsection{Fundamentos adoptados}
\label{sec:fundamentos-adoptados}

El encuadre anterior exige una arquitectura común y un modelo de dominio que
delimite y nombre lo que cada producto representa. Ambos se establecieron a
partir de los tres fundamentos de la Tabla~\ref{tab:fundamentos}, cuyas bases
conceptuales se presentaron en el marco teórico del Capítulo~\ref{cap:marco}.
Lo que aquí se fija es su adopción operativa, es decir qué se toma de cada uno,
para qué, y qué queda deliberadamente fuera.

\begin{table}[H]
\centering
{\footnotesize
\begin{tabular}{p{3.3cm}p{5.5cm}p{5.9cm}}
\toprule
\textbf{Fundamento} & \textbf{Qué aporta al marco} & \textbf{Hasta dónde se adopta} \\
\midrule
BIAN & El vocabulario de negocio y la unidad de delimitación de un producto & El Service Landscape como espacio de selección y el metadato de dominio; no gobierna el catálogo de capacidades técnicas \\
Domain-Driven Design & El modelado del interior del producto & Tres reglas del nivel estratégico y seis patrones tácticos; la relación entre \emph{bounded contexts} de proyectos distintos queda fuera \\
Arquitectura hexagonal & La separación entre el núcleo y todo lo que no es dominio & Tres principios, de los que se deriva una regla verificable por análisis estático \\
\bottomrule
\end{tabular}
}
\caption{Fundamentos adoptados y alcance de su adopción}
\label{tab:fundamentos}
\end{table}

\subsubsection*{Adopción de BIAN}

Se adopta la versión 14.0 del Service Landscape, y toda referencia a
\emph{Business Areas}, \emph{Business Domains} o \emph{Service Domains}
corresponde a esa versión. Declararla no es una formalidad, porque la
estructura del modelo cambia entre versiones y la 14.0 introdujo
renombramientos, eliminaciones e incorporaciones, particularmente en el área
de pagos.

BIAN gobierna tres aspectos del proyecto generado, que son la unidad de
derivación, el vocabulario con que se nombra lo que el producto representa y
el metadato de dominio que el producto registra. Igual de relevante es
delimitar lo que no gobierna, porque el catálogo de módulos agrupa
capacidades técnicas transversales y el estándar modela capacidades de
negocio bancario, no infraestructura técnica de aplicaciones. De los treinta
y ocho \emph{Business Domains} del Service Landscape, nueve corresponden al
tipo de productos que construye el contexto, y el subconjunto de
\emph{Service Domains} contenidos en ellos constituye el espacio de selección
por defecto del mecanismo de derivación.

\paragraph*{Unidad de derivación y criterio de granularidad}

Cada ejecución del mecanismo produce un proyecto que aloja exactamente un
\emph{bounded context}, correspondiente a un \emph{Service Domain}, y no se
admite derivar productos que alojen varios. La regla se apoya en el alcance
mínimo que define el estándar, porque un \emph{Service Domain} es una
capacidad de negocio discreta cuyo alcance es elemental, de modo que si se
descompone más, las partes pierden su contexto de negocio y pasan a ser
utilitarias. La propiedad de no solapamiento aporta la contraparte, porque
las capacidades que el estándar define no se solapan y la información que
gobiernan tampoco. Trasladada al contexto, impide que dos proyectos sean
dueños de la misma información.

\paragraph*{Registro del dominio como metadato}

Todo proyecto derivado registra, como metadato, el \emph{Service Domain} al
que corresponde y el \emph{Business Domain} que lo contiene. La regla admite
una excepción y establece una prohibición, de modo que si el producto
corresponde a una capacidad que no figura en el espacio de selección por
defecto, se permite declararla fuera de ese subconjunto registrando el
\emph{Business Domain} que la contiene. Lo que no se admite es omitir el
metadato, porque sin él la delimitación de dominio queda sin efecto
verificable y la alineación con el estándar se vuelve meramente declarativa.

\subsubsection*{Adopción de Domain-Driven Design}

Al corresponder cada proyecto a un único \emph{bounded context}, el nivel
estratégico se reduce a tres reglas, que fijan la correspondencia entre el
\emph{bounded context} y el \emph{Service Domain}, la obligación de un
lenguaje ubicuo derivado del estándar y la prohibición de alojar más de un
contexto por producto. La relación entre \emph{bounded contexts} de proyectos
distintos queda fuera del marco y se registra como punto de extensión.

En el nivel táctico se adoptan seis patrones, de los cuales cuatro son el
núcleo mínimo del modelado táctico, que son la entidad, el objeto de valor, el
agregado y el repositorio; los otros dos, el evento de dominio y el servicio
de aplicación, se incorporan por decisiones ya tomadas en el marco. Los
documentos del marco de referencia detallan cada uno y su realización en la
estructura del proyecto.

El término módulo exige una precisión, porque en este trabajo designa un
\emph{core asset} del catálogo. El patrón táctico homónimo no se adopta con ese
nombre, para evitar que un mismo término designe dos cosas distintas.

\subsubsection*{Adopción de la arquitectura hexagonal}

La arquitectura hexagonal aporta la regla que separa el núcleo del producto
de todo lo que no es dominio. De sus tres principios adoptados se deriva una
regla verificable que el marco toma como criterio, según la cual el núcleo no
importa ningún artefacto del exterior. Al construirse los productos en
TypeScript, esa regla puede comprobarse mediante análisis estático en lugar
de depender de revisión manual, lo que la convierte en una restricción
efectiva y no en una recomendación.

La articulación entre Domain-Driven Design y arquitectura hexagonal es
directa, porque los patrones tácticos describen qué contiene el núcleo, y la
arquitectura hexagonal describe cómo el núcleo se relaciona con lo que está
afuera. El repositorio es el punto donde ambos se encuentran, porque es un
patrón táctico cuya realización concreta vive en el exterior.

\subsubsection*{Aplicación conjunta de los tres fundamentos}

Los tres fundamentos operan de forma conjunta en cada derivación, y la
Tabla~\ref{tab:derivacion} recorre el camino completo, desde la selección en
el Service Landscape hasta el esqueleto generado, sobre una capacidad
concreta.

\begin{table}[H]
\centering
{\footnotesize
\begin{tabular}{p{3.6cm}p{4.6cm}p{6.5cm}}
\toprule
\textbf{Paso} & \textbf{Qué fundamento interviene} & \textbf{Qué produce} \\
\midrule
Selección de la capacidad & BIAN, Service Landscape & El \emph{Service Domain} y el \emph{Business Domain} que lo contiene \\
Delimitación del producto & BIAN y Domain-Driven Design & Un \emph{bounded context} único, con su lenguaje ubicuo \\
Modelado del núcleo & Domain-Driven Design, nivel táctico & La entidad raíz del agregado, sus objetos de valor, el contrato del repositorio y los eventos \\
Separación del exterior & Arquitectura hexagonal & Los puertos que el núcleo declara y los adaptadores que los realizan \\
Derivación del esqueleto & El mecanismo de la Fase 2 & La estructura de carpetas, los nombres y los contratos vacíos \\
\bottomrule
\end{tabular}
}
\caption{Recorrido de una derivación, desde el estándar hasta el proyecto generado}
\label{tab:derivacion}
\end{table}

El resultado sobre un producto de la familia es un núcleo en dos capas, la de
dominio dentro de la de aplicación, con los adaptadores que lo invocan desde
fuera por encima y los que él gobierna por debajo. La capacidad del estándar
no forma parte de esa arquitectura, sino que es lo que determina el contenido
del núcleo.

Lo que la derivación entrega es el punto de partida de un producto y no el
producto, porque las reglas de negocio las escribe después el equipo que lo
recibe, y esa distinción es la que mantiene al artefacto dentro de su
alcance.

\subsubsection*{Exclusión razonada de ISO 20022}

El estándar ISO 20022 se revisó como posible fundamento del marco y se
excluyó de forma deliberada. La razón es de naturaleza y no de dificultad, porque
ISO 20022 es un estándar de mensajería financiera que define modelos de
negocio y esquemas de mensajes para el intercambio de información entre
instituciones financieras, mientras que el artefacto de este trabajo no
produce, transforma ni transporta mensajes financieros, sino estructuras de
proyecto y componentes técnicos transversales. La relación entre ambos
estándares delimita el punto de contacto con precisión, porque BIAN
referencia correspondencias hacia ISO 20022 en la capa de especificación de
mensajes, y esa capa es exactamente la que el artefacto no produce.

La adopción de ISO 20022 corresponde, por tanto, a la lógica de negocio que
cada equipo implementa sobre el proyecto derivado, explícitamente fuera del
alcance de la solución. La exclusión se registra como punto de extensión y no
como omisión, porque un módulo de dominio que encapsule mensajería ISO 20022
sería un consumidor legítimo de este marco y no una parte de él.

Esta exclusión constituye una divergencia respecto del objetivo general
formulado en el anteproyecto, que comprometía la alineación con BIAN e ISO
20022. Se declara como tal, con su razón, y el Capítulo~\ref{cap:evaluacion}
la recoge entre las divergencias del proyecto al medir el cumplimiento de los
objetivos.

\section{Resumen del capítulo}

Este capítulo convirtió una intuición en un problema delimitado y un problema
en un encuadre. Al abrirlo existía la percepción de que arrancar un proyecto
cuesta demasiado y de que las funcionalidades transversales se copian; al
cerrarlo existen quince temas sostenidos por nueve entrevistas, tres
brechas que ninguna de las ocho soluciones estudiadas cubre, veintidós
requerimientos verificables y una arquitectura de la familia de productos que
la herramienta deberá materializar.

Tres decisiones estructurales sostienen el resultado, y la primera es el
encuadre como línea de producto de software, adoptado no por afinidad
conceptual sino porque los cuatro elementos que el enfoque exige ya estaban
presentes de forma informal, incluida la derivación por copia, cuya
limitación documentada es la que se encontró en campo. La segunda es la
delimitación de un producto por una capacidad de negocio del estándar, que da
a la línea una unidad de derivación y un vocabulario que no hay que inventar.
Y la tercera es que el marco diseña la familia de productos y no la
herramienta, lo que permitió verificar el encuadre antes de que existiera
código.

El capítulo dejó además dos límites declarados que el cierre del proyecto
recoge. La exclusión de ISO 20022 es una divergencia respecto del objetivo
general, y se declara con su razón en lugar de omitirse. Y las dos
desviaciones respecto del modelo de referencia clásico, el producto autónomo
y la derivación por etapas, no son concesiones sino la consecuencia de
aplicar el enfoque a un contexto de equipos autónomos con repositorios
independientes.

Lo que queda abierto por construcción es todo el mecanismo, porque el marco
acota lo que las fases siguientes pueden decidir pero no construye ni un
activo ni una línea de la herramienta. Con él quedan asignadas tres reglas
concretas que se definen donde se construyen: la forma que debe tener un
proyecto válido de la familia, los contratos por los que una capacidad
compartida se le acopla, y las condiciones que hacen válido a un módulo del
catálogo. Las dos primeras abren el Capítulo~\ref{cap:generacion}; la tercera
es del Capítulo~\ref{cap:gestor}.
